\section{Preprocessing}
\label{sec:preprocessing}

Korean is agglutinative: a single space-delimited word (\emph{eojeol}) combines a
stem with particles and endings, e.g.\ 봤는데 = 보/VV + 었/EP + 는데/EC (``saw,
but\ldots''). Treating eojeols as tokens therefore produces a very sparse vocabulary in
which the same verb appears in dozens of surface forms---a problem that is most severe
exactly in the small-data regime this study is about. All preprocessing below is
rule-based or uses a fixed, pre-trained analyser, and is applied identically to every
method; nothing is fitted on the evaluation data.

\paragraph{Normalization.}
Every review passes through six fixed rules, in order:
(1)~HTML entities are unescaped;
(2)~Unicode is normalized to \textbf{NFC}. We deliberately avoid NFKC, which maps the
compatibility jamo common in reviews (ㅋ, ㅠ, ㅎ) to a different Unicode block and can
split visually identical tokens;
(3)~zero-width and control characters are removed;
(4)~Latin letters are lowercased (Hangul is unaffected);
(5)~runs of three or more identical characters are capped at two
(ㅋㅋㅋㅋㅋ $\to$ ㅋㅋ, !!!!! $\to$ !!), which keeps the cue while merging length
variants that would otherwise be separate, rare tokens;
(6)~whitespace is collapsed. Digits, punctuation, emoji, jamo-only expressions,
English words and profanity are kept: they are frequent, informative sentiment cues.
We perform no spacing correction, no stopword removal and no digit replacement---a
review such as ``10점 만점에 1점'' (``1 point out of 10'') is legitimate content whose
digits carry the label. Cleaning (empty texts, conflicting duplicates, exact
duplicates) operates on the normalized text (\cref{tab:data}).

\paragraph{Morphological tokenization.}
The main experiment tokenizes with Kiwi~\cite{kiwipiepy} into morphemes tagged with
their part of speech (\texttt{form/TAG}, e.g.\ 좋/VA, 안/MAG), with typo correction
disabled. We keep \emph{all} morphemes, including particles and endings, because many
of them carry sentiment: negation (안/MAG, 못/MAG, -지 않-), contrast (-는데/EC) and
tone (-네/EF). Kiwi embodies external linguistic knowledge, but it is applied equally to
all thirteen methods, so it does not favour either family; the ablation in
\cref{sec:ablation} replaces it by eojeol and character tokenization to test how much
the conclusions depend on it. Sentiment symbols survive tokenization (verified
automatically): ㅋ, ㅠ, ㅎ or emoji occur in \symbolShareNSMC{} of NSMC,
\symbolShareShopping{} of Shopping and \symbolShareSteam{} of Steam reviews
(\cref{tab:a5sym}). \nKiwiEmptyNSMC, \nKiwiEmptyShopping{} and \nKiwiEmptySteam{}
reviews yielded no Kiwi tokens and were dropped.

\paragraph{Features.}
TF-IDF methods use morpheme unigrams and bigrams. sklearn's default token pattern
silently discards one-character tokens such as 좋, 안, 못, 다; we therefore pass
pre-tokenized lists, and a test asserts that 안/MAG is in the fitted vocabulary.
Neural methods use a vocabulary built on the training subset with a minimum frequency
of 2, so that the \texttt{<unk>} embedding is trained. Sequences are truncated from the
end at \maxSeqLenKiwi{} morphemes (the smallest multiple of 8 covering 95\% of the
train pool of the dataset with the longest reviews), which truncates
\truncKiwiNSMC, \truncKiwiShopping{} and \truncKiwiSteam{} of the train-pool reviews.

\paragraph{Data sparsity at small sizes.}
With 500 training reviews, \oovSmallNSMC{} of NSMC test morphemes are out of vocabulary,
falling to \oovLargeNSMC{} at 8{,}000 (Shopping: \oovSmallShopping{} $\to$ \oovLargeShopping;
Steam: \oovSmallSteam{} $\to$ \oovLargeSteam; \cref{tab:a5,fig:oov}). Neural models that
learn embeddings from scratch cannot use these tokens at all, which is one mechanism
behind their weakness at small sizes.
